\documentclass[11pt,a4paper]{article}
\usepackage[margin=1in]{geometry}
\usepackage[T1]{fontenc}
\usepackage{times}
\usepackage{authblk} 
\usepackage{amsmath} % Added for robust math support

\title{GenFSDiff: Few-Shot Steered Diffusion for Synthetic Satellite Telemetry Anomaly Generation in Support of Satellite Operations}

% Assign authors and affiliations using matching numbers
\author[1]{Joon Kim}
\author[2]{Nils-Holger Kaul}
\affil[1]{Student, Space and Ground Segment Technology, Galileo Competence Center, DLR}
\affil[2]{Research Associate, Space and Ground Segment Technology, Galileo Competence Center, DLR}

% Optional: Make the affiliation font slightly smaller and italicized
\renewcommand\Affilfont{\itshape\small}

\date{}

\begin{document}
\maketitle
Satellite telemetry anomaly detection is essential for scalable spacecraft operations but remains challenging due to the limited number of public datasets like the ESA Anomaly Detection Benchmark (ESA-ADB) and the scarcity of confirmed anomalies therein. This constraint hinders the development of high-performance health monitoring models, motivating the generation of synthetic anomalies to improve Time Series Anomaly Detection (TSAD). However, existing anomaly generation methods mostly rely on ad-hoc, hand-crafted strategies applied to raw time series. These approaches often fail to capture diverse, complex anomalous patterns and cannot generate anomalies from scratch (e.g., from pure noise). To overcome these shortcomings, this paper introduces a novel synthetic anomaly generation method: Generator for Anomalies in Satellite Telemetry using Few-shot Steered Diffusion \textbf{(GenFSDiff)}. 

GenFSDiff is trained on nominal and rare-nominal windows of six ESA-ADB Mission 1 channels (41–46). Anomaly generation can be initiated either by moderately corrupting a nominal window or directly from pure Gaussian noise. During the denoising process, the sample being generated is steered to produce realistic, diverse, univariate anomalies instead of reconstructing the original parent nominal. 

This generative steering uses a dual-objective mechanism. First, an anomaly-targeting objective pulls the generated time series towards the characteristics of a specific, known anomaly example (hence, ``few-shot''). Second, a ``shell'' constraint confines the generated sequence to the extreme boundary of the nominal data distribution (such as the 99th percentile of nominal behavior). Ultimately, this dual-force steering drives the signal toward the desired anomaly type, enforcing a clear separation from nominal behavior while penalizing implausible deviations to guarantee highly realistic outputs.

We evaluate GenFSDiff against a basic hand-crafted baseline and several state-of-the-art anomaly generation models, including the variational autoencoder-based GenIAS. Following established methodology, we measure generation quality using Anomalous Representation Proximity (ARP) for realism and the Entropy-Based Diversity Index (EDI) for diversity. Our results demonstrate that GenFSDiff achieves comparable or superior realism, nearly doubles the diversity score, and introduces a unique capability to target specific anomaly types. Furthermore, unlike existing methods, our trained model does not require raw time series as parent sequences. It can generate anomalies from pure noise without significantly degrading its ARP and EDI scores.

In the context of space operations, this framework enables the generation of realistic, diverse, and type-targeted synthetic telemetry anomalies. This capability directly supports the training and evaluation of high-performance anomaly detectors, ultimately facilitating early, automated fault detection and reducing operator workload. We will present the GenFSDiff architecture, detail our comparative evaluation on the ESA-ADB, and demonstrate how this synthetic pipeline can be deployed to enhance automated health monitoring and safeguard spacecraft operations.
\medskip

\end{document}
